% chapters/ch23_bayesian.tex
\chapter{Bayesian and Variational Reformulations}
\label{ch:bayesian}

This chapter recasts hierarchical reconstruction as a problem of Bayesian inference over tree topologies and develops the variational approximation that connects oracle-augmented algorithms to probabilistic reasoning. In the Bayesian framing, oracle responses are observations that update a prior distribution over $\mathcal{T}_n$, and the accumulated posterior concentrates on trees consistent with the majority vote of all oracle queries. The partial HC tree emerges as a natural compressed representation of this posterior: the resolved portion carries nearly all the posterior probability, while super-vertices represent regions where the posterior remains diffuse. We then introduce the variational free energy, which provides a tractable lower bound on the log-evidence and connects inference over discrete tree spaces to continuous optimization. The ELBO decomposition---expected log-likelihood minus KL divergence from the prior---reappears in Appendix~\ref{app:variational}.



\section{Bayesian Inference over Tree Metrics}

The hierarchical clustering problem has a natural Bayesian formulation. Let $P_0(T)$ be a prior distribution over tree topologies and let the observed triplet relations $\Rset = \{r_1, \ldots, r_t\}$ be conditionally independent given $T$ with likelihoods $P(r_i \mid T)$. The posterior is
\[
P(T \mid \Rset) \propto P_0(T) \prod_{r \in \Rset} P(r \mid T).
\]

Under the noisy oracle model, $P(r \mid T) = p$ if $T \models r$ and $P(r \mid T) = (1-p)/2$ otherwise. The posterior concentrates on trees consistent with the majority vote of all oracle responses.

\section{The Partial HC Tree as a Belief State}

The partial HC tree is a compressed representation of the posterior $P(T \mid \Rset)$. The structure above the super-vertex level is essentially certain (probability approaching 1 for large $n$). The uncertainty is concentrated within the super-vertices, where the local posterior is approximately uniform over the $(2|S_\mathbf{v}| - 3)!!$ local topologies.

\section{Variational Inference}

When the exact posterior is intractable to compute, variational inference approximates it with a tractable distribution $Q$.

\begin{definition}[Variational free energy]
For a variational distribution $Q$ over $\Tree_n$, the \emph{free energy} is
\[
\mathcal{F}(Q) = \mathbb{E}_Q[-\log P(\Rset \mid T)] + \KL(Q \| P_0(T)).
\]
\end{definition}

\begin{proposition}[Free energy as upper bound]
$\mathcal{F}(Q) \geq -\log P(\Rset)$ for all $Q$, with equality when $Q = P(T \mid \Rset)$.
\end{proposition}

\begin{proof}
This is a standard result in variational Bayes, following from the non-negativity of KL divergence.
\end{proof}

Minimizing $\mathcal{F}(Q)$ over a tractable family of distributions $Q$ provides a principled approximate inference procedure for tree structure.

\section{Mean-Field Approximation for Trees}

A tractable mean-field approximation assumes that the internal node assignments are independent. Under this assumption, the free energy decomposes into a sum over internal nodes, making optimization tractable. The resulting approximation is related to the partial HC tree construction: the high-confidence resolved nodes correspond to the converged mean-field assignments.

\section{Exercises}

\begin{enumerate}
  \item Write out the log-posterior $\log P(T \mid \Rset)$ as a function of the number of consistent and inconsistent triplets in $\Rset$ under the noisy oracle model.
  \item Show that minimizing free energy is equivalent to maximizing the evidence lower bound (ELBO).
  \item Describe a mean-field approximation for trees on $n$ leaves and identify which parameters the variational distribution must specify.
\end{enumerate}
